\chapter*{Предисловие}
\addcontentsline{toc}{chapter}{Предисловие}
\markboth{Предисловие}{Предисловие}

Учебник адресован школьникам 10–11 классов, изучающим информатику на
\emph{продвинутом уровне}, а также тем, кто готовится к олимпиадам,
профильным экзаменам и проектным работам в области искусственного
интеллекта. Цель книги — показать, что современный ИИ — это не
магия и не «чёрный ящик», а аккуратная математика и инженерия,
доступные старшекласснику, владеющему школьным курсом алгебры и
началами анализа.

Изложение построено по принципу \emph{от формулы к коду и обратно}.
Каждый ключевой алгоритм сопровождается:
\begin{itemize}
  \item исторической справкой и геометрической интуицией;
  \item математически строгим выводом и теоремой о свойствах;
  \item рабочей реализацией на языке Python;
  \item численным экспериментом с графиками и обсуждением границ
        применимости.
\end{itemize}

Параграфы, помеченные звёздочкой ($\star$), необязательны при первом
прочтении: в них собран более продвинутый материал, рассчитанный на
читателя, готового к доказательствам. Однако мы рекомендуем
вернуться к ним позднее: именно эти разделы дают понимание, \emph{почему}
методы машинного обучения работают, а не просто \emph{как} ими
пользоваться.

Условные обозначения, используемые в книге:

\noindent\begin{tabular}{@{}p{4cm}p{10cm}@{}}
$\R, \N, \Z$ & множества вещественных, натуральных и целых чисел \\
$f'(x),\, f''(x)$ & первая и вторая производные функции $f$ \\
$\norm{\,\cdot\,}$ & норма (длина) вектора \\
$\argmin\limits_{x} f(x)$ & точка минимума функции $f$ \\
\end{tabular}

\bigskip
\noindent\textit{Желаем интересного чтения!}
